\documentclass[UTF8]{ctexart}
\usepackage{amsmath,amssymb}
\usepackage{array}


\begin{document}
\section*{浙江大学2025--2026学年春夏学期\\线性代数II课程期末考试卷}

\begin{enumerate}
\item 已知两个平面 $\pi_1$ 为
\[
-2x+y-z-2=0,
\]
$\pi_2$ 为
\[
x+y+2z-8=0,
\]
求到 $\pi_1$ 和 $\pi_2$ 距离相等的几何平面。

\item 设 $U$ 是 $V$ 的子空间，$\pi$ 是 $V$ 到 $V/U$ 的商映射，假如 $T$ 是 $V$ 到 $W$ 的线性映射，证明存在线性映射 $S$ 从 $V/U$ 到 $W$ 满足
\[
T=S\pi
\]
充分必要条件是 $\operatorname{null}T$ 包含了 $U$。

\item 设
\[
T:\mathbb F^3\to\mathbb F^2,\qquad
T(x,y,z)=(4x+5y-6z,\;2x-3y+7z)
\]
是 $\mathbb F^3$ 到 $\mathbb F^2$ 的线性映射，
\[
\rho_i(x_1,x_2,x_3)=x_i,\qquad 1\leq i\leq 3,
\]
\[
\varphi_j(y_1,y_2)=y_j,\qquad j=1,2.
\]
求
\[
T'(\varphi_1-2\varphi_2)
\]
在 $\rho_1,\rho_2,\rho_3$ 下的坐标。

\item 设
\[
T(x,y,z)=(x+y+z,\;x+y+z,\;x+y+z)
\]
是 $\mathbb F^3$ 上的线性变换，求 $T$ 的一个二维不变子空间和 $T$ 的极小多项式。

\item 设 $V$ 是 $3$ 维实内积空间，证明
\[
\langle u+iv,\;w+ix\rangle
=
\langle u,w\rangle+\langle v,x\rangle
+i\bigl(\langle v,w\rangle-\langle u,x\rangle\bigr)
\]
是复化空间 $V_{\mathbb C}$ 上的内积。假如 $e_1,e_2,e_3$ 是 $V$ 的规范正交基，求向量组
\[
ie_1-e_2,\qquad e_1-ie_3,\qquad e_1+e_2+e_3
\]
经 Gram--Schmidt 正交化后得到的规范正交基。

\item 已知
\[
T(x,y,z)=(2y,\;3z,\;-x),
\]
求 $T$ 的奇异值和极分解。

\item $T$ 是复数域上有限维线性空间上的算子，$k$ 是正整数，证明
\[
\dim\operatorname{null}(T^k)
\leq
k\dim\operatorname{null}(T).
\]

\item 设
\[
T(x,y,z)=(3x+y-z,\;-2x-2z,\;x+y+z),
\]
在复数域中求其 Jordan 标准型，在实数域中求一组基使得在该基下对应的矩阵为
\[
\operatorname{diag}(A,A),
\qquad
A=
\begin{pmatrix}
b&-c\\
c&b
\end{pmatrix}.
\]

\item 下列论断是否正确，正确给予简单证明，不正确请举反例。

\begin{enumerate}
\item $T$ 是实线性空间 $V$ 上线性算子，$T_{\mathbb C}$、$V_{\mathbb C}$ 分别是它们的复化，那么 $V$ 是 $T_{\mathbb C}$ 不变子空间。
\item 任何复线性空间上的线性变换都存在平方根。
\item $T$ 是实内积空间 $V$ 上的线性变换，$b^2<c$，那么
\[
T^2+bT+cI
\]
是可逆的。
\item $S,T$ 是两个正算子，那么
\[
\operatorname{null}(S+T)
=
\operatorname{null}S\cap\operatorname{null}T.
\]
\end{enumerate}

\end{enumerate}

\end{document}
